\section{力矩与角动量的几何意义}

动量描述平动，角动量则同时包含运动量和它相对于参考点的分布。对固定点$O$，质点角动量为$\bm L_O=\bm r\times\bm p$，大小为$mvb$，其中$b$是$O$到速度所在直线的垂距。它不是物体必须绕$O$转圈才存在的量；一个沿直线飞过的质点，只要其运动直线不过$O$，相对$O$就有角动量。方向由右手定则确定，在平面问题中可规定逆时针为正。

力矩同样由叉积定义：$\bm\tau_O=\bm r\times\bm F$。在直角坐标中，垂直图面的分量为$\tau_z=xF_y-yF_x$。力臂是参考点到力作用线的垂距，不是参考点到作用点的距离本身。若力作用线通过$O$，即使力很大，其相对$O$的力矩仍为零。这使固定铰点和中心力问题特别适合使用角动量。

\fig{fig:em5:lever}{角动量与力矩使用相应作用线的垂距}{\begin{tikzpicture}[scale=.9]\coordinate(O)at(0,0);\coordinate(P)at(3,1.5);\fill(O)circle(.05)node[left]{$O$};\fill(P)circle(.07)node[above]{$P$};\draw[guide](O)--(P)node[midway,above]{$\bm r$};\draw[flow](P)--(5,1.5)node[above]{$\bm p$};\draw[guide](0,1.5)--(3,1.5);\draw[dim](-.4,0)--(-.4,1.5)node[midway,left]{$b$};\node at(3,-.6){$L_{Oz}=-pb$};\end{tikzpicture}}

图\ref{fig:em5:lever}中质点在$O$上方水平向右运动，$\bm r\times\bm p$指向纸里，因此逆时针为正时$L_{Oz}=-pb$。垂距$b$是$O$到动量所在直线的距离，图中的斜段$r$只在叉积中取其垂直分量。

由$\bm L_O=\bm r\times\bm p$对时间求导，第一项为$\dot{\bm r}\times m\bm v=\bm v\times m\bm v=0$，所以固定点$O$有$\dot{\bm L}_O=\bm r\times\bm F=\bm\tau_O$。对系统求和时，一对沿两质点连线的内力合力矩为$(\bm r_i-\bm r_j)\times\bm F_{ij}=0$，因此系统角动量变化只由外力矩决定。这里不仅用到了作用与反作用大小相等，还用了它们沿同一直线的性质。

对选定时间积分，得到角冲量定理$\bm L_+-\bm L_-=\int\bm\tau_{\text{外}}\dd t$。碰撞期间有限重力矩的积分可因时间很短而忽略，但固定铰点冲量的力矩是否为零取决于参考点：对铰点自身取矩为零，对其他点一般不为零。选择参考点是在消去不需要求的约束冲量，同时必须保留该点到入射作用线的垂直距离作为力臂。

\section{参考点选择与守恒分量}

角动量依赖参考点，因而“角动量守恒”必须说明相对谁守恒。固定轴承的冲量通常很大，线动量可能不守恒；但轴承冲量对轴的力矩为零，绕轴的角动量可以守恒。对绕竖直轴运动的物体，重力和竖直支持力未必各自力矩为零，但它们绕竖直轴的力矩分量为零，于是只有竖直角动量分量守恒。

\fig{fig:em5:hinge-impact}{固定铰点的冲量不改变绕铰点角动量，却通常改变线动量}{\begin{tikzpicture}

\coordinate(O)at(0,2.5);\coordinate(P)at(2,0);\draw[world](O)--(P);\fill(O)circle(.07)node[above]{$O$};\fill(P)circle(.12)node[right]{$M$};
\fill(-1.4,0)circle(.08)node[above]{$m$};\draw[flow](-1.4,.5)--(-.1,.5)node[above]{$v$};\draw[guide](-1.7,0)--(2.5,0);
\draw[dim](-2.1,0)--(-2.1,2.5)node[midway,left]{$b$};\draw[force](O)--(1.1,2.8)node[right]{$\bm J_O$};
\node[align=left]at(5.3,1.4){$\bm r_O\times\bm J_O=0$\\$L_{O,-}=mvb$\\粘合时机械能仍会损失};
\end{tikzpicture}}

图\ref{fig:em5:hinge-impact}中，铰点冲量的力臂为零，但入射动量的力臂为$b$。两者必须分别判断。若改用正在加速的任意参考点，简单的$\dot{\bm L}=\bm\tau$一般需要增加修正项；不能把一个运动接触点当作永远固定的铰点。本讲优先选固定惯性点，或者选系统质心。关于质心，外力矩等于相对质心角动量的变化率，这使自由刚体的偏心冲击可以避开平动的干扰。下一讲将给出质心角动量与刚体转动的统一表达。

平面绕定轴的质点系有$L_O=\sum m_ir_i^2\omega=I_O\omega$。这里$I_O$由质量到轴的距离决定。轻杆可不计转动惯量，但杆端的球不可不计；粘合碰撞会改变系统质量分布，故必须分别用碰前角动量与碰后转动惯量，而不能沿用同一个$I$。

\section{角动量的原点平移与质心分解}

取两个固定点$O,A$，从$O$到$A$的矢量为$\bm a$。同一质点相对两点的位置满足$\bm r_A=\bm r_O-\bm a$，所以$\bm L_A=\bm L_O-\bm a\times\bm p$。对系统求和，便有
\[\bm L_A=\bm L_O-\bm a\times\bm P.\]
因此系统总动量为零时，关于任意两个固定原点的总角动量相同；总动量不为零时，角动量数值会随原点改变。比较碰撞前后的角动量时必须固定参考点，不能碰前对一个点取矩、碰后又换到另一个点而不补上这一项。

进一步把每个质点位置写成$\bm r_i=\bm R_C+\bm\rho_i$，速度写成$\bm v_i=\bm V_C+\bm u_i$。由质心定义$\sum m_i\bm\rho_i=0$、$\sum m_i\bm u_i=0$，展开叉积求和得到
\[\bm L_O=\bm R_C\times M\bm V_C+\sum_i\bm\rho_i\times m_i\bm u_i.\]
第一项是质心绕$O$运动的角动量，第二项是相对质心的内部角动量。自由刚杆的偏心碰撞既改变第一项也改变第二项；下一讲用线冲量与绕质心角冲量分别计算，正是这一分解的直接应用。

\fig{fig:em5:origin}{位置矢量的三角形分解：先到质心，再到第i个质点}{\begin{tikzpicture}

\coordinate(O)at(0,0);\coordinate(C)at(2,.8);\coordinate(P)at(3.2,3);\fill(O)circle(.05)node[left]{$O$};\fill(C)circle(.05)node[below]{$C$};\fill(P)circle(.06)node[right]{$i$};
\draw[flow](O)--(C)node[midway,below]{$\bm R_C$};\draw[flow](C)--(P)node[midway,right]{$\bm\rho_i$};\draw[flow](O)--(P)node[midway,left]{$\bm r_i$};
\draw[flow](C)--(3.6,.8)node[right]{$\bm V_C$};
\end{tikzpicture}}

图\ref{fig:em5:origin}中从$O$到质点$i$，既可直接沿$\bm r_i$，也可先到质心$C$再沿$\bm\rho_i$。两条路径给出$\bm r_i=\bm R_C+\bm\rho_i$，正是前述角动量展开的位置分解；$\bm V_C$箭头则对应整体平动项。

如果参考点$A$本身移动，而角动量仍用各质点对惯性系的动量定义，则不能无条件沿用固定点形式。
先对 $\bm r_{iA}=\bm r_i-\bm R_A$ 求导，得 $\dot{\bm r}_{iA}=\bm v_i-\bm V_A$。
再对 $\bm L_A=\sum_i\bm r_{iA}\times m_i\bm v_i$ 使用乘积求导：
\begin{align*}
\dot{\bm L}_A
&=\sum_i(\bm v_i-\bm V_A)\times m_i\bm v_i
+\sum_i\bm r_{iA}\times m_i\dot{\bm v}_i\\
&=-\bm V_A\times\sum_i m_i\bm v_i+\bm\tau_A.
\end{align*}
其中 $\bm v_i\times m_i\bm v_i=0$，成对共线内力矩也相消，因此
\[\frac{\dd\bm L_A}{\dd t}=\bm\tau_A-\bm V_A\times\bm P.\]
附加项来自移动原点的位置变化。若$A$取系统质心，$\bm P=M\bm V_C$与$\bm V_C$平行，附加项为零，因此绕质心仍可使用$\dot{\bm L}_C=\bm\tau_C$。这解释了自由杆碰撞为何可同时对质心用线冲量、对质心轴用角冲量，而随意选一个正在加速的端点就需要格外小心。

\section{把三维角动量投影到固定轴}

\fig{em5:sphere-projection}{球面运动的子午面投影；方位速度垂直该截面}{
\begin{tikzpicture}[scale=.95]
\draw[world](0,0)circle(1.844);
\draw[axis](0,-2.1)--(0,2.8)node[above]{$z$};
\draw[guide](0,1.2)ellipse(1.4 and .28);
\draw[world](0,0)--(1.4,1.2)node[midway,below]{$R$};
\draw[guide](0,1.2)--(1.4,1.2);
\draw[dim](0,1.7)--(1.4,1.7)node[midway,above]{$R\sin\theta$};
\draw[world](0,.6)arc(90:40.6:.6);
\node at(.25,.78){$\theta$};
\fill(0,0)circle(1.5pt)node[below left]{$O$};
\fill(1.4,1.2)circle(2pt)node[right]{$P$};
\node[right]at(1.85,1.2){$\otimes\ v_\phi$};
\draw[flow,blue](1.4,1.2)--(2.2,.27)node[right]{$v_\theta$};
\node at(0,-2.65){$\phi$沿纬圈计，$\theta$沿子午线计};
\end{tikzpicture}}

图\ref{em5:sphere-projection}中，关于$z$轴的力臂是水平距离$R\sin\theta$，方位速度为$v_\phi=R\sin\theta\,\dot\phi$；子午速度$v_\theta=R\dot\theta$在截面内，对$L_z$没有贡献。因此
\[
 L_z=m(R\sin\theta)v_\phi
     =mR^2\sin^2\theta\,\dot\phi.
\]
图中$\otimes$表示所选位置的正方位速度向纸面里。极角的转折只令$v_\theta=0$，若$L_z\ne0$，方位运动仍然存在。下面讨论守恒条件时，始终取同一固定轴。


\fig{fig:em5:axis-torque}{悬点到摆球的水平距离决定竖直角动量，重力矩本身却沿水平方向}{\begin{tikzpicture}

\coordinate(O)at(0,3);\coordinate(P)at(2,0);\draw[world](O)--(P);\draw[axis](0,-.7)--(0,3.8)node[above]{$z$};\draw[guide](0,0)ellipse(2 and .35);
\fill(O)circle(.05)node[left]{$O$};\fill(P)circle(.09)node[above right]{$P$};\draw[force](P)--(2,-1.2)node[below]{$mg$};
\node at(2.8,.4){$\otimes\ v_\phi$};\node[align=left]at(5.4,2){$\tau_z=xF_y-yF_x=0$\\$L_z=mr_\perp v_\phi$};
\draw[dim](0,-.65)--(2,-.65)node[midway,below]{$r_\perp$};
\end{tikzpicture}}

三维运动中，力矩为零的某个分量就足以产生相应角动量分量的守恒。取固定竖直轴$z$，有$\dot L_z=\tau_z$。重力沿竖直方向，所以它关于任意点的力矩虽未必为零，其竖直分量必为零。若其他约束力也没有绕该轴的力矩，$L_z$便守恒。这正是球碗中方位运动、自由转台以及可转动螺旋导线的共同依据。

由图\ref{em5:sphere-projection}的轴向角动量表达式可见，在非零$L_z$保持不变时，小球靠近竖直轴会使方位速度增大，形成阻止继续靠轴的等效势垒。极角的转折点只是$\dot\theta=0$，不能把该点误设成完全静止。

角动量关于点的完整矢量，与关于轴的分量不可混用。图\ref{fig:em5:axis-torque}中的圆锥摆，关于悬点的角动量有不断旋转的水平分量，重力力矩恰好维持它的方向变化；其竖直分量才守恒。若改对轨道圆心取矩，合力始终径向，完整的竖直角动量反而可以保持不变。这些结论不矛盾，因为参考点不同，所取分量也不同。

\section{中心力、面积速度与径向变化}

\fig{em5:polar-sector}{径向、切向速度与半径扫过的微小面积}{
\begin{tikzpicture}[scale=1.05]
\fill[blue!10](0,0)--(15:2.4)arc(15:32:2.4)--cycle;
\draw[world](0,0)--(15:3.3) (0,0)--(32:3.3);
\draw[world](15:2.4)arc(15:32:2.4);
\fill(0,0)circle(1.5pt)node[below left]{$O$};
\coordinate(P)at(15:2.4);
\fill(P)circle(2pt)node[below]{$P$};
\draw[flow,blue](P)--($(P)+(.97,.26)$)node[right]{$v_r=\dot r$};
\draw[flow,red](P)--($(P)+(-.31,1.16)$)node[above]{$v_\theta=r\dot\theta$};
\node[below]at(15:1.3){$r$};
\draw[world](15:.8)arc(15:32:.8);
\node[right]at(24:1){$\dd\theta$};
\node at(24:1.8){$\dd A$};
\end{tikzpicture}}

图\ref{em5:polar-sector}把速度分为沿半径与垂直半径的两部分。径向分量不产生关于$O$的角动量，也不扫出一阶面积；微小扇形面积为$\dd A=\tfrac12r^2\dd\theta$。所以$L=mr v_\theta$与$\dot A=L/(2m)$来自同一几何关系。半径本身变化时也可在当前$r$处作这一瞬时分解，不能把合速度直接当作切向速度。


中心力始终指向或背离固定点$O$，故对$O$力矩为零，图\ref{em5:polar-sector}中的角动量保持不变。半径变化时，总速率$v=\sqrt{\dot r^2+v_\theta^2}$仍包含径向运动，不能直接代入$L=mrv$；只有纯切向运动的瞬间才可这样写。小孔绳长变化及三维碗内纬圈运动，都需要把径向、经向或方位速度分开。

图中单位时间扫过的面积保持不变，角速度却未必不变；当半径缩小时，$\dot\theta=L/(mr^2)$增大，切向速率$L/(mr)$也增大。所需增加的动能可以来自拉绳者的功或重力势能转化，角动量守恒本身并不能给出能量是否增加。

结合能量，可以把角向运动写成径向的等效势能$U_{\mathrm{eff}}(r)=U(r)+L^2/(2mr^2)$。第二项不是新增的一种真实势能，而是把守恒角动量对应的切向动能放到只研究$r$的方程中。径向转折点满足$\dot r=0$，但切向速度一般不为零，不能把“半径最大”误写成“小球瞬时静止”。

\section{中心力圆轨道的稳定性}

中心吸引力$F_r=-kr^{-n}$（$k>0$）提供一个进一步检验角动量与能量配合的方法。由于$L=mr^2\dot\theta$守恒，径向方程为$m\ddot r=-kr^{-n}+L^2/(mr^3)$。圆轨道$r=r_0$要求两项平衡，$L^2/(mr_0^3)=kr_0^{-n}$。有圆轨道并不保证它稳定，还需看小径向偏离后的净力方向。

令 $r=r_0+\xi$ 且 $|\xi|\ll r_0$，记净径向力
$f(r)=-kr^{-n}+L^2r^{-3}/m$。在固定 $L$ 下逐项求导，
\[
f'(r)=nk r^{-n-1}-\frac{3L^2}{m}r^{-4}.
\]
圆轨道条件给 $L^2/m=kr_0^{3-n}$，所以
$f'(r_0)=(n-3)kr_0^{-n-1}$。
一阶展开 $f(r_0+\xi)\simeq f(r_0)+f'(r_0)\xi$ 中第一项为零，得到
\[m\ddot\xi=-(3-n)k r_0^{-n-1}\xi.\]
若 $n<3$，令 $\omega_r^2=(3-n)k/(mr_0^{n+1})$，方程为 $\ddot\xi+\omega_r^2\xi=0$。
试解 $\e^{st}$ 给根 $s=\pm i\omega_r$，故
$\xi=A\cos\omega_rt+B\sin\omega_rt$。
求导得 $\dot\xi=-A\omega_r\sin\omega_rt+B\omega_r\cos\omega_rt$；
给定 $\xi(0)=\xi_0,\dot\xi(0)=w_0$，得到
$A=\xi_0,B=w_0/\omega_r$，偏离有界并作小振动。

若 $n>3$，令 $\kappa^2=(n-3)k/(mr_0^{n+1})$，特征方程变为 $s^2-\kappa^2=0$，通解为
$\xi=A\e^{\kappa t}+B\e^{-\kappa t}$。
初值给 $A+B=\xi_0,\kappa(A-B)=w_0$，所以
$A=(\xi_0+w_0/\kappa)/2$、$B=(\xi_0-w_0/\kappa)/2$。
一般微扰含增长项，圆轨道不稳定。

$n=3$ 时圆轨道条件为 $L^2=mk$，代回完整径向方程，净力对所有 $r$ 都为零。此时 $\ddot r=0$，以 $\dot r(0)=w_0,r(0)=r_0+\xi_0$ 积分：
$\dot r-w_0=\int_0^t0\,\dd\tau=0$，
$r-(r_0+\xi_0)=\int_0^tw_0\dd\tau=w_0t$。
纯位置偏移没有回复力，径向初速度会导致持续漂移；这不是稳定的小振动。这里固定的是小扰动过程中守恒的 $L$，不是强令角速度不变。

\fig{em5:radial-force}{固定角动量时，圆轨道附近净径向力的符号；横坐标$\rho=r/r_0$}{
\begin{tikzpicture}[x=2.8cm,y=.9cm]
\begin{scope}
\draw[axis](.55,0)--(1.85,0)node[right]{$\rho$};
\draw[axis](1,-2.3)--(1,1.9)node[above]{$F_{\rm eff}/(kr_0^{-n})$};
\draw[world,blue,domain=.65:1.7,samples=80]plot(\x,{-1/(\x)^2+1/(\x)^3});
\node[below left]at(1,0){$1$};
\node at(1.1,-2.9){$n=2<3$：回复};
\end{scope}
\begin{scope}[xshift=5.7cm]
\draw[axis](.55,0)--(1.85,0)node[right]{$\rho$};
\draw[axis](1,-2.3)--(1,1.9)node[above]{$F_{\rm eff}/(kr_0^{-n})$};
\draw[world,red,domain=.65:1.7,samples=80]plot(\x,{-1/(\x)^4+1/(\x)^3});
\node[below left]at(1,0){$1$};
\node at(1.1,-2.9){$n=4>3$：偏离增大};
\end{scope}
\end{tikzpicture}}

图\ref{em5:radial-force}中正力指向半径增大方向。左图在$\rho>1$时力为负、在$\rho<1$时力为正，均指回圆轨道；右图的力与偏离方向相同。两图在$r_0$处的斜率正是展开式中$-(3-n)kr_0^{-n-1}$的符号。分析沿同一守恒角动量进行，因此等效的离心项也随$r$改变，不能固定角速度作此比较。


圆轨道角频率满足$\Omega^2=k/(mr_0^{n+1})$，因此$\omega_r=\sqrt{3-n}\,\Omega$。平方反比吸引力$n=2$使两频率相同；这个结果来自具体力律与稳定性展开，面积速度守恒本身并不能单独确定轨道形状或径向振动频率。这一频率关系还用到了具体力律，不能只凭面积速度守恒推出。

\section{多个守恒量的配合}

复杂题中，线动量、角动量和机械能分别承担不同任务。线动量约束质心运动，角动量处理力矩和偏心作用，机械能联系速度大小、势能和耗散。是否使用某条方程取决于所选系统的受力与约束，而不取决于题目被归入哪个专题。

图\ref{fig:em5:two-stages}把刚粘合与碰后上摆分开：冲击阶段的 $\omega_+$ 成为摆动阶段的初始角速度，随后质心上升 $\Delta h$ 才进入重力势能变化。碰撞与连续运动常须分开：固定轴处的粘合冲击可用绕轴角动量守恒，但不能用机械能守恒；碰后摆动若轴无摩擦，机械能又可守恒。自由车上的发射，水平方向线动量守恒、竖直方向一般不守恒；发射器储能转化为球和车两部分动能。先在图上标出参考系中的速度，再用几何关系把相对速度转成对地速度，能避免重复或遗漏动能。

\fig{fig:em5:two-stages}{绕固定轴的粘合冲击与随后的摆动使用不同的能量条件}{\begin{tikzpicture}

\foreach \xx/\ang/\lab in {0/0/刚粘合,4.5/55/向上摆动}{\begin{scope}[xshift=\xx cm]
\fill(0,2.5)circle(.06)node[above]{$O$};\draw[world](0,2.5)--({2*sin(\ang)},{2.5-2*cos(\ang)});\fill({2*sin(\ang)},{2.5-2*cos(\ang)})circle(.12);
\draw[guide](0,2.5)--(0,.5);\node at(.6,3.2){\lab};\end{scope}}
\draw[flow](0,.5)--(1.2,.5)node[above]{$\omega_+\ell$};\draw[dim](7.2,.5)--(7.2,1.353)node[midway,right]{$\Delta h$};
\node at(1,-.4){冲击：$L_{O,-}=I_+\omega_+$};\node at(5.7,-.4){摆动：$K_+=K+\Delta U$};
\end{tikzpicture}}

\section{活动约束中的功与角动量交换}

“光滑约束力不做功”应理解为对固定约束、相对位移沿其切线而言。活动的导线、槽或曲面可以把能量从一个物体传给另一个，即使接触处没有摩擦。可转动螺旋导线上的小环，在相对导线的速度方向上不受约束力分量，但导线自身正在转动，约束力可以分别改变小环和导线的动能。把两者合为系统，接触作用的总功抵消，机械能才可整体守恒。

\fig{fig:em5:helix}{相对螺旋转角决定下降高度，导线反向转动使绝对方位角速度减小}{\begin{tikzpicture}

\draw[axis](0,-.4)--(0,5.5)node[above]{$z$};\draw[guide](0,4.6)ellipse(1.2 and .25);
\draw[world,domain=0:720,samples=180]plot({1.2*cos(\x)},{4.6-\x/180+.25*sin(\x)});
\draw[dim](0,5.05)--(1.2,5.05)node[midway,above]{$R$};\fill(.6,2.716827)circle(.1)node[left]{$m$};
\draw[flow](.6,2.716827)--(1.38,2.572)node[right]{相对滑动};
\draw[flow](-1.4,.2)..controls(-.5,-.2)and(.5,-.2)..(1.4,.2);\node at(0,-.7){导线转角$\psi$，惯量$I$};
\draw[dim](2,2.6)--(2,4.6)node[midway,right]{$2\pi c$};\node[align=left]at(5,3.4){相对转角$\phi$增大\\$z_0-z=c\phi$\\对地转角$\psi+\phi$};
\end{tikzpicture}}

图\ref{fig:em5:helix}给出具体模型。螺旋半径$R$、螺距$2\pi c$，可绕固定竖直轴无摩擦转动，转动惯量$I$；小环质量$m$。令导线转角为$\psi$，小环相对导线沿下降方向的转角为$\phi$。几何给出$z=z_0-c\phi$，小环对地角速度为$\dot\psi+\dot\phi$。系统从静止释放，绕轴无外力矩，故
\[
 I\dot\psi+mR^2(\dot\psi+\dot\phi)=0,\qquad
 \dot\psi=-\frac{mR^2}{I+mR^2}\dot\phi.
\]
记 $\eta=I/(I+mR^2)$，则 $\dot\psi=-(1-\eta)\dot\phi$，
小环对地角速度为 $\dot\psi+\dot\phi=\eta\dot\phi$，竖直速度为 $\dot z=-c\dot\phi$。
先将三部分动能写全：
\[
K=\frac12I\dot\psi^2
+\frac12mR^2(\dot\psi+\dot\phi)^2+\frac12mc^2\dot\phi^2.
\]
角向两项的系数为
$I(1-\eta)^2+mR^2\eta^2
=ImR^2/(I+mR^2)=mR^2\eta$，
所以 $K=m(c^2+\eta R^2)\dot\phi^2/2$。
下降 $h$ 时，由机械能守恒得到
\[
 \tfrac12m(c^2+\eta R^2)\dot\phi^2=mgh.
\]
这时下降速率$c\dot\phi$比导线固定时更大，但小环的对地总速率平方为$(c^2+\eta^2R^2)\dot\phi^2<2gh$，因为导线也分得一部分动能。因此比较“速度”必须指明是对地总速率、竖直分量还是相对导线速率。

\fig{fig:em5:sleeve}{固定套管提供水平外力，杆的允许速度却是竖直方向}{\begin{tikzpicture}

\draw[world](0,0)--(0,3.6);\draw[world](-.35,1.3)--(-.35,2.3) (.35,1.3)--(.35,2.3);\draw[world](-1,1.3)--(-.35,1.3) (-1,2.3)--(-.35,2.3);
\node[left]at(-1,1.8){固定套管};\draw[world](0,0)--(2.5,0);\draw[world](1.5,1)arc(180:360:1);\fill(3.3,.4)circle(.08);
\draw[force](0,1.8)--(1.5,1.8)node[right]{$N_h$};\draw[flow](.8,3.5)--(.8,2.6)node[right]{$V_z$};
\node[align=left]at(5.5,1.3){$W_{N_h}=0$\\$\Delta P_x=\int N_h\dd t$一般不为零};
\end{tikzpicture}}

图\ref{fig:em5:sleeve}中固定套管限制竖直杆，情况又不同。套管对杆、槽系统的水平约束力可产生外冲量，使水平动量不守恒；但杆的受力点沿竖直方向移动，水平约束力的功率为零，仍不输入机械能。因此动量和机械能应分别判断，不能把“有外力”概括成所有守恒量都不可用。

\fig{fv:e5:support}{活动支点的等长约束限制相对速度，张力功率仍取对地速度}{\begin{tikzpicture}

\coordinate(O)at(0,1.8);\coordinate(P)at(2,-.5);\draw[world](O)--(P);\fill(O)circle(.05)node[above]{$O$};\fill(P)circle(.075)node[below]{$P$};
\draw[force](P)--(1.1,.535)node[midway,left]{$\bm T$};\draw[flow](O)--(1.6,1.8)node[right]{$\bm V_O$};
\draw[flow](P)--(4.175,0)node[midway,above]{$\bm v_P$};\draw[guide](P)--(3.6,-.5)node[midway,below]{$\bm V_O$};\draw[flow](3.6,-.5)--(4.175,0)node[midway,right]{$\bm v_{\text{相}}$};
\node at(2,-1.5){$\bm v_{\text{相}}\perp OP,\quad\bm v_P=\bm V_O+\bm v_{\text{相}}$};
\end{tikzpicture}}
图\ref{fv:e5:support}中绳长不变只给 $(\bm v_P-\bm V_O)\cdot\bm n=0$，因此张力对质点的功率为 $\bm T\cdot\bm v_P=\bm T\cdot\bm V_O$。支点移动时它一般非零，不能继续按固定支点的圆周运动假定张力不做功。同时 $O$既然在运动，就应使用活动参考点的角动量变化关系；“力的作用线通过 $O$”本身不足以保证关于该活动点的角动量守恒。


\section{分离后用新的动力学衔接}

细绳不能推，未粘连的接触面不能拉。因此守恒方程只给候选状态，还须检查张力、支持力是否非负。弹簧托住但未连接的物块，压缩量到零且相对运动仍向外时应结束接触；此后物块与托盘的加速度不同，不能继续当作一个整体。

设分离瞬间物块相对托盘向上的速度为$w>0$，以此刻接触高度差为相对坐标零点。若随后物块自由下落、托盘的向上加速度为常量$a_T$，则图\ref{fig:em5:relative-flight}的相对坐标为
\[
 \zeta(t)=wt-\tfrac12(g+a_T)t^2.
\]
只有$g+a_T>0$，在这一阶段内才有第二个相遇根$t=2w/(g+a_T)$；还须确认托盘在此之前没有改变受力或碰到边界。若$a_T$不是常量，应积分其实际运动，不能套固定抛出点的$2v/g$。

\fig{fig:em5:relative-flight}{分离后间距先增后减，相遇时间取相对位移的第二个零点}{\begin{tikzpicture}
\draw[world](-.7,0)--(1.2,0);\fill(.2,1.4)circle(.1)node[right]{物块};\draw[force](.2,1.4)--(.2,.5)node[left]{$mg$};\draw[flow](1.4,-.7)--(1.4,.4)node[right]{$a_T$正向向上};
\begin{scope}[xshift=5.3cm]\draw[axis](0,0)--(4.7,0)node[right]{$t$};\draw[axis](0,0)--(0,2.4)node[above]{$\zeta$};
\draw[blue,thick,domain=0:2,samples=80]plot({2*\x},{1.8*(2*\x-\x*\x)});\node[below]at(4,0){$2w/(g+a_T)$};\node at(2,2.5){相对间距};\end{scope}\end{tikzpicture}}

衔接时保留三组信息：分离的位置、各物体对地速度、各自下一阶段的受力。能量给出的可能是系统速度组合，须先用约束关系分配给各物体；相遇后若发生冲击，再换用新的冲量和恢复条件。这些中间状态共同决定后续过程，不能只保留一个最终速度数值。
